% =============================================================================
%  书末：部分习题答案与提示
% =============================================================================
\unnumberedchapter{部分习题答案与提示}

\begingroup
\small
\setlength\parindent{0pt}
\setlength\parskip{.35\baselineskip}

\subsection*{第 1 章}

{\mojingheadfont\color{mjAccent} 1.}\quad
(1)~$x=1$，$y=2$，$z=3$。\quad
(2)~第三个方程等于第二个减去第一个，有无穷多解：$x_1=3-t$，$x_2=2+t$，$x_3=t$（$t\in\R$）。\quad
(3)~消元后得 $0=-1$，无解。

{\mojingheadfont\color{mjAccent} 2.}\quad
$\begin{bmatrix}1&0&1\\0&1&1\\0&0&0\end{bmatrix}$。

{\mojingheadfont\color{mjAccent} 3.}\quad
系数行列式为 $(\lambda+2)(\lambda-1)^2$。$\lambda\ne1$ 且 $\lambda\ne-2$ 时有唯一解
$x=y=z=\frac{1}{\lambda+2}$；$\lambda=1$ 时三个方程都化为 $x+y+z=1$，有无穷多解；$\lambda=-2$ 时
三式相加得 $0=3$，无解。

{\mojingheadfont\color{mjAccent} 4.}\quad
设鸡翁、母、雏各 $x$，$y$，$z$ 只，则 $x+y+z=100$，$5x+3y+\frac{z}{3}=100$。消去 $z$ 得
$7x+4y=100$，故 $4\mid x$。令 $x=4t$，则 $y=25-7t$，$z=75+3t$，$t=0,1,2,3$，即
$(0,25,75)$，$(4,18,78)$，$(8,11,81)$，$(12,4,84)$。《张丘建算经》给出的是后三组。

{\mojingheadfont\color{mjAccent} 7.}\quad
若 $\vb{u}$，$\vb{v}$ 是解，则 $A\bigl(t\vb{u}+(1-t)\vb{v}\bigr)=t\vb{b}+(1-t)\vb{b}=\vb{b}$；
$\vb{u}\ne\vb{v}$ 时，不同的 $t$ 给出不同的解。

\subsection*{第 2 章}

{\mojingheadfont\color{mjAccent} 1.}\quad $10$。

{\mojingheadfont\color{mjAccent} 2.}\quad $(a+3b)(a-b)^3$。提示：把各列加到第一列，提出 $a+3b$。

{\mojingheadfont\color{mjAccent} 3.}\quad $\det A=-7$，$\det A_1=-14$，$\det A_2=-7$，故 $x=2$，$y=1$。

{\mojingheadfont\color{mjAccent} 5.}\quad
$\det A=\det A\T=\det(-A)=(-1)^n\det A=-\det A$。

{\mojingheadfont\color{mjAccent} 7.}\quad
若 $\det A=\pm1$，由克拉默法则，解的分量都是整数。反之，取 $\vb{b}$ 为单位向量 $\vb{e}_j$，
得整数矩阵 $X$ 使 $AX=I_n$，于是 $\det A\cdot\det X=1$，两者都是整数，故 $\det A=\pm1$。
\endgroup
